% a05_twoloop_v2: proof of lem:a05:D45 (v), part 2: derivative bounds, box inequality, polarisation remainder, the unit-mass bound, conclusion (input from a05_cnum_NEG.tex)
\begin{lemma}[derivative bounds at unit mass]\label{lem:a05:EB}
Let $a\le1$ and $q,k\in\R^3$.
\begin{enumerate}[label=\textup{(\roman*)},nosep,leftmargin=*]
\item $\big|\sqrt{\hat e(q+k)}-\sqrt{\hat e(q)}\big|\le|k|$; hence $C(q+k)\le2C(q)$ and $\chi(q+k)\le2\chi(q)$ for $|k|\le\frac12$.
\item $a^2\le13/\chi(q)$.
\item $|\partial_\mu\hat e|\le2\sqrt{\hat e}$, and every derivative of $\hat e$ of order $\ge2$ is bounded by $2$. Consequently, for $|k|\le\frac12$ and all multi-indices $b,\beta$: $|\partial^b_q\partial^\beta_kC(q+k)|\le C_{b,\beta}\,\chi(q)^{-1-(|b|+|\beta|)/2}$, $|\partial^b_q\partial^\beta_k\tilde V_\mu(q,q+k)|\le C_{b,\beta}\,\chi(q)^{(1-|b|-|\beta|)/2}$, and the derivatives of $\cos\frac{a(2q+k)_\mu}2$ and of the phases $e^{\pm\mathrm iak_\mu/2}$ are bounded by $C_{b,\beta}\,\chi(q)^{-(|b|+|\beta|)/2}$: each derivative gains a factor $\chi(q)^{-1/2}$ relative to the bound of the undifferentiated factor ($C(q+k)\le2\chi(q)^{-1}$, $|\tilde V_\mu(q,q+k)|\le2\sqrt{\hat e(q)}+|k|\le3\chi(q)^{1/2}$, $|\cos|,|e^{\mathrm i\cdot}|\le1$).
\item $|P_\rho(k)|\le C\sqrt{\hat e(k)}$ for all $k$, and $|P_\rho(k)-k_\rho|\le Ca|k|^2$ for $|k|\le\frac12$.
\end{enumerate}
\end{lemma}
\begin{proof}
(i) With $v_\mu(q)\defeq\frac2a\sin\frac{aq_\mu}2$, $\sqrt{\hat e}=|v|$ and $|v_\mu(q+k)-v_\mu(q)|\le|k_\mu|$ by $|\sin x-\sin y|\le|x-y|$; the triangle inequality gives the first claim. If $x\defeq\sqrt{\hat e(q)}\ge\frac12$, then $\hat e(q+k)+1\ge(x-\frac12)^2+1\ge\frac{x^2+1}2$ (equivalent to $(x-1)^2+\frac12\ge0$); if $x<\frac12$, then $\hat e(q+k)+1\ge1\ge\frac{x^2+1}2$. Also $\hat e(q+k)+1\le(x+\frac12)^2+1\le2x^2+2$.
(ii) $\hat e\le12/a^2$ and $a\le1$.
(iii) $\partial_\mu\hat e=\frac2a\sin(aq_\mu)=\frac4a\sin\frac{aq_\mu}2\cos\frac{aq_\mu}2$, of modulus $\le2|v_\mu|$; $\partial_\mu\partial_\nu\hat e=2\cos(aq_\mu)\delta_{\mu\nu}$, and each further derivative multiplies by $a\le1$. A derivative of $C=\chi^{-1}$ produces $C^2\partial\hat e$ with $|\partial\hat e|\le2\chi^{1/2}$, a gain of $2\chi^{-1/2}$; a derivative of a factor $\partial\hat e$ gives a factor $\le2\le2\chi^{1/2}\chi^{-1/2}$. The $j$-th derivative of $\tilde V_\mu(q,q+k)=-\frac2a\sin\frac{a(2q+k)_\mu}2$ is bounded by $2a^{j-1}\le2\cdot13^{(j-1)/2}\chi(q)^{-(j-1)/2}$ for $j\ge1$ by (ii), and that of $\cos\frac{a(2q+k)_\mu}2$ or $e^{\pm\mathrm iak_\mu/2}$ by $a^j\le13^{j/2}\chi(q)^{-j/2}$. For the factor $C(q+k)$ and its derivatives use $\chi(q+k)\ge\frac12\chi(q)$ from (i). Induction on $|b|+|\beta|$ gives the stated bounds.
(iv) Applying \cite[(3-2)]{Rei88b} to $F\equiv q_\rho$ gives $P_\rho(0)=0$ and $\nabla P_\rho(0)=e_\rho$, so $p_\rho(x)=x_\rho+O(|x|^2)$, and by periodicity $p_\rho$ vanishes on $2\pi\Z^3$. Hence $|p_\rho(x)|\le\|\nabla p_\rho\|_\infty\dist(x,2\pi\Z^3)\le\|\nabla p_\rho\|_\infty\frac\pi2\big(4\sum_\mu\sin^2\frac{x_\mu}2\big)^{1/2}$ by \cref{lem:a05:el}(c), and $|P_\rho(k)|=a^{-1}|p_\rho(ak)|\le C\sqrt{\hat e(k)}$; and $|P_\rho(k)-k_\rho|=a^{-1}|p_\rho(ak)-ak_\rho|\le a^{-1}C|ak|^2=Ca|k|^2$ for $|ak|\le\frac12$.
\end{proof}

\begin{lemma}[box inequality]\label{lem:a05:box}
\textup{(i)} For $g\in C^d$ on a cube $Q_d\subset\R^d$ of unit side, $\sup_{Q_d}|g|\le\sum_{S\subset\{1,\dots,d\}}\int_{Q_d}|\partial_Sg|$, $\partial_S\defeq\prod_{i\in S}\partial_i$ ($\partial_\emptyset g=g$). \textup{(ii)} For $G$ smooth and $\frac{2\pi}a$-periodic in $(k_1,k_2)\in\R^6$,
\[
\sup_{k_2\in Q}\big\langle|G(\cdot,k_2)|\big\rangle_{a,k_1}\le\sum_{S\subset\{1,2,3\}}\int_Q\big\langle|\partial^{k_2}_SG(\cdot,k_2)|\big\rangle_{a,k_1}\,dk_2,\qquad
\sup_{Q^2}|G|\le\sum_{S\subset\{1,\dots,6\}}\int_{Q^2}|\partial_SG| .
\]
\end{lemma}
\begin{proof}
(i) For $x,y\in Q_d$ the fundamental theorem of calculus, applied coordinate by coordinate, gives $g(x)=\sum_S\int_{[y_S,x_S]}(\partial_Sg)(t_S,y_{S^c})\,dt_S$, so $|g(x)|\le\sum_S\int_{[0,1]^{|S|}}|(\partial_Sg)(t_S,y_{S^c})|\,dt_S$ for every $y$; average over $y\in Q_d$. (ii) Apply (i) in $k_2$ at fixed $k_1$ and integrate over $k_1$ (Fubini), respectively apply (i) in the six variables.
\end{proof}
Applied to $G=\chi(k_2)R_{\mathrm{dbl}}F^{(1)}$, respectively $G=\chi(k_1)\chi(k_2)R_{\mathrm{dbl}}F^{(1)}$: on $Q$, $\chi\le\frac74$, $|\partial_i\chi|=|\frac2a\sin(ak_i)|\le2|k_i|\le1$, and the mixed derivatives of $\chi$ vanish ($\hat e$ is a sum of functions of one coordinate each), so by the Leibniz rule both right sides of (ii) are $\le CK$, with $K$ of \cref{lem:a05:H} (derivatives of order $\le3$, respectively $\le6$) and a numerical $C$.

\begin{lemma}[the polarisation remainder]\label{lem:a05:VPR}
Let $F_b(q;k)\defeq e^{-\mathrm ia(k_\mu-k_\nu)/2}\tilde V_\mu(q,q+k)\tilde V_\nu(q+k,q)C(q)C(q+k)$ be the bubble integrand at unit mass with photon indices $\mu,\nu$, or either of its doubled versions ($C(q)\to C(q)^2$ or $C(q+k)\to C(q+k)^2$). For $a\le1$ and $|k|\le\frac12$,
\[
\big|\big\langle(1-t^1_k)F_b(\cdot;k)\big\rangle_a\big|\le C_{\mathrm{VP}}|k|^2,
\]
and the left side vanishes at $k=0$. Moreover $\Pi_{\mathrm{dbl}}(0)\defeq\langle F_{b,\mathrm{dbl}}(\cdot;0)\rangle_a$ obeys $|\Pi_{\mathrm{dbl}}(0)|\le4\langle C^2\rangle_a\le\frac{\pi^2}{16}$ and $\nabla_k\langle F_{b,\mathrm{dbl}}(\cdot;k)\rangle_a\big|_{k=0}=-\frac{\mathrm ia}2(e_\mu-e_\nu)\Pi_{\mathrm{dbl}}(0)$, of modulus $\le a|\Pi_{\mathrm{dbl}}(0)|$.
\end{lemma}
\begin{proof}
For fixed $q$, $(1-t^1_k)F_b(q;k)=\big[F_b(q;k)-F_b(q;0)-k\cdot\nabla_kF_b(q;0)\big]+\big(k-P(k)\big)\cdot\nabla_kF_b(q;0)$. The bracket equals $\int_0^1(1-\theta)\,k^{\mathsf T}\nabla_k^2F_b(q;\theta k)\,k\,d\theta$, and by \cref{lem:a05:EB}(iii), $|\partial^\beta_kF_b(q;k')|\le C\chi(q)^{-1-|\beta|/2}$ for $|k'|\le\frac12$ (two vertex factors of size $\chi^{1/2}$, two propagators $\chi^{-1}$, each derivative gaining $\chi^{-1/2}$; the doubled versions are smaller by a factor $2\chi^{-1}$). Hence $\langle|\nabla^2_kF_b(\cdot;\theta k)|\rangle_a\le C\langle\chi^{-2}\rangle_a\le C$ by \eqref{eq:a05:soft}, and the bracket contributes $\le C|k|^2$. For the second term, $|k-P(k)|\le Ca|k|^2$ by \cref{lem:a05:EB}(iv) and $\langle|\nabla_kF_b(\cdot;0)|\rangle_a\le C\langle\chi^{-3/2}\rangle_a\le C(1+\ln\frac1a)$ by \eqref{eq:a05:soft}; and $a(1+\ln\frac1a)\le1$ for $a\le1$. At $k=0$ the remainder vanishes identically in $q$, since $P(0)=0$. For the last claims: at $k=0$ the integrand is $\tilde V_\mu(q,q)\tilde V_\nu(q,q)C(q)^3$ with $|\tilde V_\mu(q,q)|=\frac2a|\sin(aq_\mu)|\le2\sqrt{\hat e(q)}$, so $|\Pi_{\mathrm{dbl}}(0)|\le4\langle\hat eC^3\rangle_a\le4\langle C^2\rangle_a$; and $\langle F_{b,\mathrm{dbl}}(\cdot;k)\rangle_a=e^{-\mathrm ia(k_\mu-k_\nu)/2}T_{\mathrm{dbl}}(k)$ with $T_{\mathrm{dbl}}$ even in $k$ (the substitution $q\mapsto-q$ maps the integrand at $k$ to the integrand at $-k$, the doubled line to itself), so $\nabla T_{\mathrm{dbl}}(0)=0$.
\end{proof}

\begin{proposition}[the unit-mass bound, uniform in the photon mass]\label{prop:a05:U}
With $a_0$ of \cref{lem:a05:H} there is $C_U<\infty$, depending on $a_0$ only, such that $|Y(\nu,a)|\le C_U$ for all $\nu>0$ and $a\in(0,a_0]$.
\end{proposition}
\begin{proof}
By \cref{lem:a05:forest} at $s=1$, $Y(\nu,a)$ is a finite sum of terms $c_\Gamma\langle\langle R'F^{(\nu)}_{\Gamma,\mathrm{dbl}(l)}\rangle\rangle_a$. Each is bounded by a constant of the stated kind.

\emph{(0) Vanishing terms.} $\Gamma_{11}$ has no scalar line. For $\Gamma_6$ and $\Gamma_9$ the polarisation subgraph is the seagull loop; the seagull vertex depends on its two photon momenta only through their sum, which is $0$ at zero external momentum, so the loop integrand $F_{\mathrm{VP}}(q;k)$ does not depend on $k$ and $(1-t^1_k)F_{\mathrm{VP}}\equiv0$. Hence $R'F_{\mathrm{dbl}}\equiv0$ for every $l$: if $l$ lies on the path, $R'=R_{\mathrm{dbl}}$ subtracts the loop; if $l$ lies in the loop, $R'F_{\mathrm{dbl}}=(1-t^1_k)F_{\mathrm{dbl}}=0$.

\emph{(1) Class A ($\Gamma_0,\Gamma_2,\Gamma_3,\Gamma_5,\Gamma_7,\Gamma_8,\Gamma_{10}$), the $R_{\mathrm{dbl}}$ part.} By \cref{lem:a05:PF}, $|R_{\mathrm{dbl}}F^{(\nu)}|=D_\nu(k_1)D_\nu(k_2)|R_{\Gamma,l}(k_1,k_2)|=D_\nu(k_1)\chi(k_1)\,D_\nu(k_2)\chi(k_2)\,|R_{\mathrm{dbl}}F^{(1)}|$. Split $\BZ_a^2$ according to whether $k_1$, $k_2$ are hard or soft. On hard momenta $D_\nu\chi\le1+\pi^2$. Hence
\begin{align*}
\langle\langle\mathbf 1_{hh}|R_{\mathrm{dbl}}F^{(\nu)}|\rangle\rangle_a&\le(1+\pi^2)^2\langle\langle|R_{\mathrm{dbl}}F^{(1)}|\rangle\rangle_a\le(1+\pi^2)^2K,\\
\langle\langle\mathbf 1_{hs}|R_{\mathrm{dbl}}F^{(\nu)}|\rangle\rangle_a&\le(1+\pi^2)\,\sigma_1\sup_{k_2\in Q}\big\langle|\chi(k_2)R_{\mathrm{dbl}}F^{(1)}(\cdot,k_2)|\big\rangle_{a,k_1}\le(1+\pi^2)\sigma_1CK,
\end{align*}
by \cref{lem:a05:box}(ii) and \cref{lem:a05:H}, and the same for $sh$ with $k_1\leftrightarrow k_2$; and $\langle\langle\mathbf 1_{ss}|R_{\mathrm{dbl}}F^{(\nu)}|\rangle\rangle_a\le\sigma_1^2\sup_{Q^2}|\chi(k_1)\chi(k_2)R_{\mathrm{dbl}}F^{(1)}|\le\sigma_1^2CK$. With $\sigma_1\le\frac1{16}$ from \eqref{eq:a05:soft}, $|\langle\langle R_{\mathrm{dbl}}F^{(\nu)}\rangle\rangle_a|\le C$.

\emph{(2) Class A, the removed terms of \eqref{eq:a05:OS}.} Type (d) ($\Gamma_2$, $\Gamma_5$, $\Gamma_7$): the coefficient $\langle F_{\gamma,\mathrm{dbl}}(\cdot;0)\rangle_a$ is the one-loop vertex subgraph with the line $l$ doubled at zero external momenta, and it vanishes by $q\mapsto-q$: for the triangle the integrand is $\sum_\alpha\tilde V_\alpha(0,q)\tilde V_\mu(q,q)\tilde V_\alpha(q,0)\,D_\nu(q)C(q)^3$ up to the two phases of the internal photon, which cancel, with $\tilde V_\mu(q,q)=-\frac2a\sin(aq_\mu)$ odd and the other factors even; for the swordfish vertex (the seagull with one external photon leg, its other leg joined by the internal photon to a current vertex) the integrand is $(1+e^{-\mathrm iaq_\mu})(e^{\mathrm iaq_\mu}-1)D_\nu(q)C(q)^2=2\mathrm i\sin(aq_\mu)D_\nu(q)C(q)^2$, odd. ($\Gamma_8$ has no removed term, its $A^3$-tadpole containing no scalar line; the subtraction $t^0$ of that tadpole is $0$ anyway, as $V_3\propto\sin\frac{a(p+p')_\mu}2$.) Type (a) ($\Gamma_0$ with $l$ the inner line): $t^1_\gamma F_{\gamma,\mathrm{dbl}}=F_{\gamma,\mathrm{dbl}}(k_2;0)+P(k_1)\cdot\nabla F_{\gamma,\mathrm{dbl}}(k_2;0)$, and the order-$1$ term is odd in the inner loop momentum $k_2$ (as in the proof of \cref{lem:a05:forest}(a)), so the removed term's double integral is $N_\gamma\,I_{\bar\Gamma}$ with
\[
N_\gamma=\big\langle D_\nu(k)\,\hat e(k)\,C(k)^2\big\rangle_a,\qquad I_{\bar\Gamma}=\big\langle D_\nu(k)\,\hat e(k)\,C(k)^2\big\rangle_a,
\]
the inner exchange graph with its scalar line doubled at zero external momentum, and the outer exchange graph with the contracted subgraph as a unit vertex between its two scalar lines of equal momentum; here $\sum_\mu\tilde V_\mu(0,-k)\tilde V_\mu(-k,0)=\hat e(k)$. Since $D_\nu\hat e\le1$, both are $\le\langle C^2\rangle_a\le\frac{\pi^2}{64}$.

\emph{(3) Class B ($\Gamma_1$, $\Gamma_4$).} One photon chord $k$, one polarisation chord $q$. By \cref{lem:a05:PF} every term is $D_\nu(k)^2$ times a $\nu$-free function, which is the path factor times the subtracted bubble: the path factor is $\tilde V_\mu(0,-k)\tilde V_\nu(-k,0)C(k)$ for $\Gamma_1$ (with $C(k)^2$ when $l$ is the path line), bounded entrywise by $\hat e(k)C(k)\le1$, and a constant for $\Gamma_4$ (the seagull at zero scalar momenta and zero total photon momentum). \emph{Soft $k$:} use $R'F_{\mathrm{dbl}}=D_\nu(k)^2\,\mathrm{path}(k)\,(1-t^1_k)F_{b(,\mathrm{dbl})}(q;k)$ itself ($R'$ subtracts the bubble by definition; for $l$ on the path, $R'=R_{\mathrm{dbl}}$). By \cref{lem:a05:VPR} its $q$-integral is bounded by $C_{\mathrm{VP}}|k|^2|\mathrm{path}(k)|$, and $\langle\mathbf 1_{Q_s}D_\nu(k)^2|k|^2\rangle_a\le\frac{\pi^2}{64}$ by \eqref{eq:a05:soft}, uniformly in $\nu$. \emph{Hard $k$:} $R'F_{\mathrm{dbl}}=R_{\mathrm{dbl}}F_{\mathrm{dbl}}$ minus the removed term, which is present only when $l$ lies in the bubble. The $R_{\mathrm{dbl}}$ part is $\le(1+\pi^2)^2\langle\langle|R_{\mathrm{dbl}}F^{(1)}|\rangle\rangle_a\le(1+\pi^2)^2K$. The removed part is
\[
\big\langle\mathbf 1_{\mathrm{hard}}\,D_\nu(k)^2\,|\mathrm{path}(k)|\,\big|\Pi_{\mathrm{dbl}}(0)+P(k)\cdot\nabla_k\langle F_{b,\mathrm{dbl}}(\cdot;k)\rangle_a\big|_{k=0}\big|\big\rangle_{a,k}
\le(1+\pi^2)^2\,|\Pi_{\mathrm{dbl}}(0)|\,\big\langle D_1^2\big(1+Ca\sqrt{\hat e}\big)\big\rangle_a\le C\Big(1+a\big(1+\ln\tfrac1a\big)\Big)\le C,
\]
by \cref{lem:a05:VPR}, \cref{lem:a05:EB}(iv) and \eqref{eq:a05:soft} ($D_1^2\sqrt{\hat e}\le\chi^{-3/2}$). ($\Gamma_1$ vanishes identically on $\eps\Z^3$ by the transversality noted in \cref{lem:a05:40}(b); the bound covers it regardless.)

Summing (0)--(3) over the finitely many pairs $(\Gamma,l)$ gives $|Y(\nu,a)|\le C_U$, where $C_U$ collects $K$ and numerical constants, and depends on $a_0$ only.
\end{proof}

\begin{proof}[Proof of (v)]
By \cref{lem:a05:SC}, $y_{40}(s;\eps)=s^{-1}Y(\lmu/\sqrt s,\eps\sqrt s)$. For $\eps\le a_0/\sqrt s$ the second argument lies in $(0,a_0]$, and \cref{prop:a05:U} gives $|y_{40}(s;\eps)|\le C_U/s$ for every $\lmu>0$. So (v) holds with $C_{40}\defeq C_U$, which depends neither on $\lmu$ nor on $\eps$ nor on $s$. In scaled form, $s\,y_{40}$ is bounded by a function of $\lmu/\sqrt s$ and $\eps\sqrt s$ that stays bounded as $\lmu/\sqrt s\to0$: no $\ln(\sqrt s/\lmu)$ occurs.
\end{proof}
